\documentclass[a4paper, twocolumn]{article}

\usepackage[utf8]{inputenc}
\usepackage[T1]{fontenc}
\usepackage[slovene]{babel}
\usepackage{amsmath}
\usepackage{amssymb}
\usepackage{graphicx}
\usepackage{geometry}
\usepackage{hyperref}
\usepackage{microtype}
\usepackage{float}
\usepackage{caption}
\captionsetup{skip=2pt}
\raggedbottom

\geometry{margin=2cm, top=2.5cm}

\title{\textbf{Ploščina zanke Bézierjeve krivulje}}
\author{Sinja Kočica\\
\small Univerza v Ljubljani, Fakulteta za računalništvo in informatiko}
\date{}

\begin{document}

\maketitle

\section{Opis naloge}

Izračunamo ploščino zanke, ki jo omejuje Bézierjeva krivulja dana s kontrolnim poligonom osmih točk:
\[
(0,0),\,(1,1),\,(2,3),\,(1,4),\,(0,4),\,(-1,3),\,(0,1),\,(1,0).
\]
Krivulja stopnje 7 (določena z 8 kontrolnimi točkami) tvori zanko, ki omejuje zaprto območje v ravnini. Naloga zahteva natančen izračun ploščine te zanke.

\section{Opis rešitve}

Rešitev poteka v treh korakih: evalvacija krivulje, iskanje presečišča in izračun ploščine z numerično integracijo.

\subsection{Evalvacija krivulje in De Casteljaujev algoritem}

Bézierjeva krivulja stopnje $n$ je parametrična krivulja določena z $n+1$ kontrolnimi točkami $P_0, P_1, \ldots, P_n$:

\[
\mathbf{B}(t) = \sum_{k=0}^{n} \binom{n}{k} t^k (1-t)^{n-k} P_k,
\quad t \in [0, 1].
\]

Vrednost krivulje pri danem $t$ računamo z \emph{de Casteljaujevim algoritmom}, ki rekurzivno interpolira med kontrolnimi točkami:

\[
P_k^{(r)}(t) = (1-t)\, P_k^{(r-1)}(t) + t\, P_{k+1}^{(r-1)}(t).
\]

Začnemo z $P_k^{(0)} = P_k$ in po $n$ korakih dobimo $\mathbf{B}(t) = P_0^{(n)}(t)$. Algoritem je numerično stabilnejši od neposrednega izračuna z Bernsteinovimi polinomi.

Ker je Bézierjeva krivulja polinom, lahko njen odvod izračunamo analitično. Odvod je spet Bézierjeva krivulja stopnje $n-1$ z novimi kontrolnimi točkami:

\[
P'_k = n\,(P_{k+1} - P_k).
\]

S tem se izognemo numeričnim diferencam in ohranimo strojno natančnost.

\subsection{Iskanje presečišča}

Iščemo parametra $t_1 \neq t_2$ za katera velja $\mathbf{B}(t_1) = \mathbf{B}(t_2)$. To prevedemo na sistem dveh enačb v dveh nezankah:

\[
F(t_1, t_2) = \mathbf{B}(t_1) - \mathbf{B}(t_2) = \mathbf{0}.
\]

Postopek poteka v dveh korakih. Najprej krivuljo vzorčimo na 200 točkah in poiščemo par, ki si je najbližji. To nam da začetni približek. 

Sistem $F(t_1, t_2) = 0$ rešimo z Newtonovo metodo. Na vsakem koraku izračunamo kako spremeniti $t_1$ in $t_2$, da se $F$ približa nič. Pri tem uporabimo odvode krivulje v obeh točkah.

\subsection{Izračun ploščine z Greenovo formulo}

Ko poznamo $t_1$ in $t_2$, ploščino zanke izračunamo z \emph{Greenovo formulo}:

\[
P = \frac{1}{2} \left|\int_{t_1}^{t_2}
\bigl(x(t)\,\dot{y}(t) - \dot{x}(t)\,y(t)\bigr)\,dt\right|.
\]

Namesto integriranja po celotnem 2D območju zanke integriramo le vzdolž krivulje, kar je bistveno preprosteje. Integral računamo z \emph{Gauss-Legendrovo kvadraturo} s 5 točkami na vsakem podintervalu. Z $m$ točkami je kvadratura točna za polinome stopnje do $2m-1$. V našem primeru 5 točk zadosti za krivuljo stopnje 9.

\section{Primeri uporabe}
 
\subsection{Evalvacija krivulje}
 
De Casteljaujev algoritem deluje za Bézierjeve krivulje poljubne stopnje. Spodnja slika prikazuje linearno, kvadratno in kubično krivuljo skupaj s kontrolnimi poligoni. 

Za kvadratno krivuljo z kontrolnimi točkami $(0,0)$, $(1,2)$, $(2,0)$ algoritem vrne $B(0.5) = (1, 1)$, kar sovpada z analitičnim izračunom.
 
\begin{figure}[H]
    \centering
    \includegraphics[width=\linewidth]{demo/bezier_evalvacija.png}
    \caption{Evalvacija Bézierevih krivulj stopnje 1, 2 in 3.}
\end{figure}
 
\subsection{Iskanje zanke}
 
Algoritem za iskanje zanke prav tako deluje na splošnih Bézierevih krivuljah. Za kvadratno krivuljo pravilno ugotovi, da zanka ne obstaja in vrne vrednost \texttt{nothing}. Za krivuljo, ki seka samo sebe, pa najde presečišče z natančnostjo $2.3 \times 10^{-16}$.
 
\begin{figure}[H]
    \centering
    \includegraphics[width=\linewidth]{demo/bezier_kvadrat.png}
    \caption{Kvadratna krivulja brez zanke.}
\end{figure}
 
\begin{figure}[H]
    \centering
    \includegraphics[width=\linewidth]{demo/bezier_petica.png}
    \caption{Krivulja z zanko. }
\end{figure}
 
\subsection{Krivulja iz naloge}
 
Krivulja s kontrolnim poligonom iz naloge se seče pri parametrih $t_1 \approx 0.0751$ in $t_2 \approx 0.9249$, kar odraža simetrijo kontrolnega poligona ($t_1 + t_2 \approx 1$). Točka presečišča je pri koordinatah $(0.500,\, 0.616)$.

Ploščina zanke je:

\[
P \approx 2.2537095301723.
\]

Numerično stabilnost metode preverimo s konvergenčnim testom. Ploščino izračunamo z različnim številom podintervalov in opazujemo kako se rezultat spreminja. Ko povečevanje števila podintervalov ne spremeni več rezultata, smo dosegli konvergenco. Vrednost se stabilizira na 13 decimalkah že pri 100 podintervalih.
 
\begin{figure}[H]
    \centering
    \includegraphics[width=\linewidth]{demo/bezier_naloga.png}
    \caption{Bézierjeva krivulja iz naloge. }
\end{figure}

 
\end{document}
